\documentclass[11pt,a4paper]{article}

\usepackage[utf8]{inputenc}
\usepackage[T1]{fontenc}
\usepackage[ngerman]{babel}
\usepackage[a4paper,margin=2.0cm]{geometry}
\usepackage{booktabs}
\usepackage{tabularx}
\usepackage{array}
\usepackage{amsmath}
\usepackage{hyperref}
\usepackage{xcolor}

\hypersetup{
  colorlinks=true,
  linkcolor=blue,
  urlcolor=blue,
  pdftitle={Parameterbereiche für Supernova-Trainingsdaten},
  pdfauthor={Astro-Ai-Lab}
}

\newcolumntype{Y}{>{\raggedright\arraybackslash}X}

\title{Parameterbereiche für Supernova-Trainingsdaten}
\author{Astro-Ai-Lab}
\date{9. November 2026}

\begin{document}

\maketitle

\section*{Zweck}

Dieses Dokument trennt die derzeit verwendete Referenzkonfiguration von den
vorgeschlagenen Bereichen für einen späteren Trainingsdatensatz. Die Bereiche
sind noch keine endgültige Nutzerentscheidung. Sie müssen vor der Erzeugung
größerer Datensätze mit der Betreuung abgestimmt werden.

\section{Aktuelle Referenzkonfiguration}

Die bisher geplotteten und verglichenen Supernova-Läufe basieren auf einem
gemeinsamen turbulenten Medium. Die beiden erfolgreichen Vergleichsläufe
unterschieden sich nur durch die CFL-Zahl.

\begin{tabularx}{\textwidth}{p{4.1cm}Y}
\toprule
Parameter & Aktueller Wert \\
\midrule
Auflösung & $128^3$ Zellen \\
Physikalische Boxgröße & $100\,\mathrm{pc}$ \\
Referenz-Wasserstoffdichte & $n_{\mathrm H,0}=1\,\mathrm{cm}^{-3}$ \\
Referenztemperatur & $T_0=10^4\,\mathrm K$ \\
Mittlere Molekülmasse & $\mu\approx0{,}61$; ionisiertes Gas \\
Supernovaenergie & $E_{\mathrm{SN}}=10^{51}\,\mathrm{erg}$ \\
Explosionsradius & $0{,}03$ Codeeinheiten $\approx3\,\mathrm{pc}$ \\
Glättung & $2$ Zellen \\
Forcierungswellenzahl & $k_f=25$ \\
Forcierungskorrelation & $t_{\mathrm{corr}}=1$ Codezeiteinheit \\
Forcing-Amplitude & $9$; nicht direkt eine Machzahl \\
Zufallsseed & $42$ \\
Cooling & Schure-2009-CIE, Temperaturboden $10^4\,\mathrm K$ \\
Cooling-Integration & Phase 1 explizit, Phase 2 implizit \\
Supernova-Solver & PP-WENO mit PP-Flux-Blending, REDISTRIBUTE und $\texttt{vacuum\_rest=True}$ \\
CFL-Vergleich & $0{,}05$ und $0{,}1$ erfolgreich; $0{,}5$ wird aktuell geprüft \\
\bottomrule
\end{tabularx}

\section{Vorläufige Literaturbereiche}

\begin{tabularx}{\textwidth}{p{3.5cm}p{4.0cm}Y}
\toprule
Parameter & Literaturanker & Vorläufige Auswahl für Trainingsdaten \\
\midrule
Mittlere Dichte $n_{\mathrm H}$ & Kim und Ostriker untersuchen $0{,}1$ bis $100\,\mathrm{cm}^{-3}$; Zhang und Chevalier verwenden $1$, $10$ und $100\,\mathrm{cm}^{-3}$. & $0{,}1$, $1$, $10$, $100\,\mathrm{cm}^{-3}$ \\
\addlinespace
Turbulente Machzahl $\mathcal M_s$ & Warmes ionisiertes Medium: $\mathcal M_s\lesssim2$; turbulente Molekülwolken: $\mathcal M=3$, $10$, $30$. & $1$, $2$, $3$, $10$, $30$; Regime getrennt kennzeichnen \\
\addlinespace
Supernovaenergie $E_{\mathrm{SN}}$ & Für eine einzelne Supernova ist $10^{51}\,\mathrm{erg}$ der kanonische Referenzwert. & Zunächst $0{,}5$, $1$, $2\times10^{51}\,\mathrm{erg}$ als bewusste Explorationsspanne; dies ist keine direkte vollständige Beobachtungsverteilung. \\
\addlinespace
Anfangstemperatur & Aktuelles Modell beschreibt warmes, ionisiertes Gas. & Zunächst $10^4\,\mathrm K$ konstant; kalte Molekülwolken separat behandeln. \\
\addlinespace
Cooling und Metallizität & Schure et al. liefern eine optisch dünne CIE-Kühlkurve auf SPEX-Basis. & Zunächst Schure2009-CIE mit solarer Referenzzusammensetzung und $T_{\mathrm{floor}}=10^4\,\mathrm K$. \\
\bottomrule
\end{tabularx}

\section{Wichtige Einschränkungen}

\begin{itemize}
  \item Die Forcing-Amplitude ist kein direkter Machzahl-Parameter. Für jeden
    Zielwert von $\mathcal M_s$ muss die Amplitude zunächst kalibriert werden,
    indem die tatsächliche RMS- beziehungsweise lokale Machzahl im Medium bei
    $t=5$ bestimmt wird.
  \item Die Werte $\mathcal M_s=3$, $10$ und $30$ stammen aus Modellierungen
    turbulenter Molekülwolken. Unser aktuelles Schure-Setup ist dagegen auf
    ionisiertes Gas mit $\mu\approx0{,}61$ und einem Temperaturboden von
    $10^4\,\mathrm K$ festgelegt. Für kalte Molekülwolken wären eine andere
    Skalierung, ein anderer Temperaturboden und eine passende Kühlbehandlung
    notwendig.
  \item Die räumliche Auflösung muss mit der Explosionsgröße und der
    Schalenbildungsradius verglichen werden. Kim und Ostriker empfehlen für
    konvergierte Größen näherungsweise $\Delta x,r_{\mathrm{init}}<r_{\mathrm{sf}}/3$.
    Bei $128^3$ Zellen in einer $100\,\mathrm{pc}$-Box gilt
    $\Delta x\approx0{,}78\,\mathrm{pc}$. Bei hohen Dichten kann der aktuelle
    Explosionsradius von ungefähr $3\,\mathrm{pc}$ daher zu groß sein.
  \item CFL, Solver, Flux-Blending und Positivitätsschutz sind numerische
    Parameter. Sie sollten für numerische Vergleichsläufe variiert, aber nicht
    als physikalische Trainingsparameter interpretiert werden.
\end{itemize}

\section{Empfohlene erste Staffelung}

Für einen ersten überschaubaren Datensatz wird empfohlen, zunächst das aktuelle
warme ionisierte Regime zu verwenden:

\[
n_{\mathrm H}\in\{0{,}1,1,10,100\}\,\mathrm{cm}^{-3},\qquad
\mathcal M_s\in\{1,2\},\qquad
E_{\mathrm{SN}}\in\{0{,}5,1,2\}\times10^{51}\,\mathrm{erg}.
\]

Danach kann ein getrennt gekennzeichneter Molekülwolken-Datensatz mit
$\mathcal M_s=3$, $10$ und $30$ ergänzt werden, sobald die physikalische
Skalierung und das Cooling für tiefe Temperaturen erweitert wurden.

\section{Literatur}

\begin{enumerate}
  \item C.-G. Kim und E. C. Ostriker (2015): \emph{Momentum Injection by
    Supernovae in the Interstellar Medium}, ApJ 802, 99.
    \url{https://arxiv.org/abs/1410.1537}
  \item D. Zhang und R. A. Chevalier (2018): \emph{Numerical simulations of
    supernova remnants in turbulent molecular clouds}, MNRAS 477, 4521.
    \url{https://arxiv.org/abs/1807.06603}
  \item B. M. Gaensler et al. (2011): \emph{Low--Mach-number turbulence in
    interstellar gas revealed by radio polarization gradients}, Nature 478,
    214--217. \url{https://arxiv.org/abs/1110.2896}
  \item K. M. Schure et al. (2009): \emph{A new radiative cooling curve based
    on an up to date plasma emission code}, A\&A 508, 751.
    \url{https://arxiv.org/abs/0909.5204}
\end{enumerate}

\end{document}
